% from main.tex
%%%%%%%%%%  scan idx 396  |  folio 383  %%%%%%%%%%
\[
\begin{tikzcd}[column sep=large,row sep=large]
D_+(st) \arrow[r,hook,"i"] \arrow[dr,"f_{U,st}"'] &
D_+(s) \arrow[d,"f_{U,s}"]\\
& U
\end{tikzcd}
\]
il s'%
agit de montrer que les Modules
\[
  i^*\circ\bar f_{U,s}^{\,*}\bigl(M_{(s),U}\bigr)
  \quad\text{et}\quad
  \bar f_{U,st}^{\,*}\bigl(M_{(st),U}\bigr)
\]
sont canoniquement isomorphes. Or on a d'une part les isomorphismes canoniques :
\[
\begin{aligned}
i^*\circ\bar f_{U,s}^{\,*}\bigl(M_{(s),U}\bigr)
&\simeq
i^*\Bigl(f_{U,s}^{-1}\bigl(M_{(s),U}\bigr)
 \otimes_{f_{U,s}^{-1}((B|_U)_{(s)})}\mathcal O_{D_+(s)}\Bigr)\\
&\simeq
i^{-1}\Bigl(i_{U,s}^{-1}\bigl(M_{(s),U}\bigr)
 \otimes_{f_{U,s}^{-1}((B|_U)_{(s)})}\mathcal O_{D_+(s)}\Bigr)
 \otimes_{i^{-1}(\mathcal O_{D_+(s)})}\mathcal O_{D_+(st)}\\
&\simeq
f_{U,st}^{-1}\bigl(M_{(s),U}\bigr)
 \otimes_{f_{U,st}^{-1}((B|_U)_{(s)})}\mathcal O_{D_+(st)} .
\end{aligned}
\]
D'autre part :
\[
  \bar f_{U,st}^{\,*}\bigl(M_{(st),U}\bigr)
  \simeq
  f_{U,st}^{-1}\bigl(M_{(st),U}\bigr)
  \otimes_{f_{U,st}^{-1}((B|_U)_{(st)})}\mathcal O_{D_+(st)} .
\]
Enfin, il résulte de EGA II 2.2.2, traduit en termes de faisceaux, qu'on a les isomorphismes canoniques :
\[
  M_{(st),U}\simeq
  \bigl(M_{(s),U}\bigr)_{t^d/s^e}
  \simeq
  M_{(s),U}\otimes_{(B|_U)_{(s)}}(B|_U)_{(st)},
\]
d'où finalement l'isomorphisme cherché.

Il est clair que si \(M\) est quasi-cohérent le Module \(\Proj_0(M)\) ainsi construit est bien celui qui est défini dans EGA II.

\textbf{2.3.3.} Il résulte de 2.3.1 et 2.3.2 que le foncteur \(\Proj_0\) est exact. Revenant aux notations de 2.2, on en déduit donc par passage aux catégories dérivées un foncteur encore noté \(\Proj_0\) :
\[
  \Proj_0:D^+(\mathcal C)\longrightarrow D^+(\mathcal A).
\]
On voit également sur les constructions précédentes que le
